\section{Deployment Strategy}

The cluster is deployed using Ansible playbooks that automate the provisioning of K3s across all nodes. A single control-plane node is initialized as the K3s server, after which worker nodes are bootstrapped as agents and automatically joined to the cluster via a shared token. The Ansible playbooks handle host-level configuration such as package dependencies, kernel module loading (NVIDIA drivers, container runtime), user setup, and networking, as well as the K3s installation itself. This semi-manual approach keeps the bootstrapping process transparent and auditable while reducing the risk of configuration drift across nodes.

Once the K3s cluster is operational, all platform components: JupyterHub, Volcano, Kube-Ray, MLflow, Keycloak, Harbor, and the observability stack are installed and managed exclusively through Helm charts with environment-specific value overrides. Chart lifecycle management (installation, upgrades, rollbacks) is automated via GitHub Actions pipelines, which lint, template, and apply the Helm releases on each push to the main branch. This CI/CD-driven approach ensures that every change to the cluster configuration is version-controlled, reviewable, and reproducible, supporting the single-administrator operability goal: the platform can be reconstructed from the repository on a fresh K3s installation with no manual configuration steps beyond the initial Ansible bootstrap.
